\documentclass[12pt,reqno]{amsart}

\usepackage{amsmath,amssymb,amsthm}
\usepackage{booktabs}
\usepackage[colorlinks=true,linkcolor=blue,citecolor=blue,urlcolor=blue]{hyperref}
\usepackage[margin=1.15in]{geometry}
\usepackage{enumitem}
\usepackage{microtype}

\newtheorem{theorem}{Theorem}[section]
\newtheorem{proposition}[theorem]{Proposition}
\newtheorem{corollary}[theorem]{Corollary}
\newtheorem{conjecture}[theorem]{Conjecture}
\theoremstyle{definition}
\newtheorem{definition}[theorem]{Definition}
\theoremstyle{remark}
\newtheorem{remark}[theorem]{Remark}
\newtheorem{observation}[theorem]{Observation}

\newcommand{\ZZ}{\mathbb{Z}}
\newcommand{\QQ}{\mathbb{Q}}
\newcommand{\RR}{\mathbb{R}}
\newcommand{\CC}{\mathbb{C}}
\newcommand{\Gstar}{G^{\!*}}
\newcommand{\abs}[1]{\lvert #1 \rvert}
\newcommand{\Aut}{\operatorname{Aut}}
\newcommand{\re}{\operatorname{Re}}

\title{The Ratio and the Product:\\[2pt]
  Two Objects from the Euler Reflection Formula at $z=\tfrac14$}

\author{William J.\ Steinmetz III}
\address{Independent Researcher}
\date{April 3, 2026}

\begin{document}

\begin{abstract}
The Euler reflection formula $\Gamma(\tfrac14)\,\Gamma(\tfrac34)=\pi\sqrt{2}$ constrains
but does not determine the pair $\{\Gamma(\tfrac14),\,\Gamma(\tfrac34)\}$. Two independent
objects encode the residual freedom: the \emph{product}
$\Gamma(\tfrac14)\,\Gamma(\tfrac34)=\pi\sqrt{2}$, which is algebraic in~$\pi$; and the
\emph{ratio} $\Gstar=\Gamma(\tfrac14)/\Gamma(\tfrac34)=2.958\,675\ldots$, which is
algebraically independent of~$\pi$. We show that $\log\Gstar$ absorbs every unsolved
Dirichlet $L$-value at integer argument with rational coefficients, while $\pi$ carries
every solved one. The master quadratic $x^2-16\Gstar^2\,x+16\Gstar^3=0$ built from the
ratio has roots $x_+=137.036$ and $x_-=3.024$. Three centuries of mathematics studied
the product. We argue that the ratio is where the coupling constants live, and that the
two operations---product and ratio---correspond to the two modes of information processing
that physics calls measurement and superposition.
\end{abstract}

\maketitle

%==========================================================================
\section{Two Objects from One Formula}
%==========================================================================

The Euler reflection formula at $z=\tfrac14$ reads
\begin{equation}\label{eq:reflection}
  \Gamma\!\left(\tfrac14\right)\,\Gamma\!\left(\tfrac34\right)
  \;=\; \frac{\pi}{\sin(\pi/4)}
  \;=\; \pi\sqrt{2}\,.
\end{equation}
This constrains one combination of $\Gamma(\tfrac14)$ and $\Gamma(\tfrac34)$ but leaves
one degree of freedom. That degree of freedom can be parameterized as:
\begin{alignat}{2}
  \text{The product:}&\quad \Gamma\!\left(\tfrac14\right)\,\Gamma\!\left(\tfrac34\right)
  &&= \pi\sqrt{2}\,,
  \label{eq:product}\\[4pt]
  \text{The ratio:}&\quad
  \Gstar \;:=\; \frac{\Gamma(\tfrac14)}{\Gamma(\tfrac34)}
  &&= 2.958\,675\,119\,188\,639\ldots
  \label{eq:ratio}
\end{alignat}
The product is fixed by~\eqref{eq:reflection}. The ratio $\Gstar$ is algebraically
independent of~$\pi$ \cite{Chudnovsky1980,Nesterenko1996}. Together they determine both
Gamma values individually:
\begin{equation}
  \Gamma\!\left(\tfrac14\right) = \sqrt{\Gstar\,\pi\sqrt{2}}\,,\qquad
  \Gamma\!\left(\tfrac34\right) = \sqrt{\frac{\pi\sqrt{2}}{\Gstar}}\,.
\end{equation}
The product and the ratio are the two independent pieces of information in the
pair $\{\Gamma(\tfrac14),\,\Gamma(\tfrac34)\}$.

%==========================================================================
\section{What the Product Contains}
%==========================================================================

The product $\pi\sqrt{2}$ lives in the world of solved constants.
The anti-correlation theorem \cite{Steinmetz2026a} establishes that at each integer
$s\geq 2$, exactly one of $\zeta(s)$ and $\beta(s)$ reduces to a rational multiple
of a power of~$\pi$. The solved values are:
\begin{itemize}[nosep]
  \item Even $\zeta$: $\;\zeta(2)=\pi^2/6$, $\;\zeta(4)=\pi^4/90$, \ldots
  \item Odd $\beta$: $\;\beta(1)=\pi/4$, $\;\beta(3)=\pi^3/32$, \ldots
\end{itemize}
These are the $L$-values that reduce to~$\pi$. Physics built on~$\pi$ obtained Gaussian
integration, statistical mechanics, the path integral measure---but not the coupling
constants.

%==========================================================================
\section{What the Ratio Contains}
%==========================================================================

The ratio $\Gstar$ absorbs the complementary set: the unsolved $L$-values.

\begin{theorem}[Log $\Gstar$ identity]\label{thm:loggstar}
\begin{equation}\label{eq:loggstar}
  \log\Gstar \;=\; \frac{\gamma+3\log 2}{2}
  \;-\;\frac{G}{2}
  \;+\;\frac{7}{24}\,\zeta(3)
  \;-\;\frac{\beta(4)}{4}
  \;+\;\frac{31}{160}\,\zeta(5)
  \;-\;\cdots
\end{equation}
where $G=\beta(2)$ is Catalan's constant, $\zeta(3)$ is Ap\'ery's constant, and the
general coefficients are: $\zeta(2m+1)$ has coefficient
$(2^{2m+1}-1)/\bigl((2m+1)\cdot 2^{2m+1}\bigr)$; $\beta(2m)$ has coefficient
$(-1)^{m+1}/(2m)$.
\end{theorem}

Every unsolved constant appears: Catalan, Ap\'ery, $\beta(4)$, $\zeta(5)$, and the
entire infinite tower. PSLQ analysis at 80+ digits finds no algebraic relation of
degree~$\leq 20$ between any individual unsolved constant and
$\{\pi,\Gstar,\varpi\}$.

The master quadratic built from $\Gstar$ is
\begin{equation}\label{eq:quadratic}
  x^2 - 16\Gstar^2\,x + 16\Gstar^3 = 0\,,
\end{equation}
with coefficient $16=\abs{\Aut(E_i)}^2$ where $E_i\colon y^2=x^3-x$ is the CM elliptic
curve with $j=1728$. The roots are
\begin{equation}
  x_+ = 137.036\,171\ldots\,,\qquad x_- = 3.023\,964\ldots
\end{equation}
The ratio contains the fine structure constant. The product does not.

%==========================================================================
\section{Product as Collapse, Ratio as Distinction}
%==========================================================================

The two operations act differently on information:

\begin{observation}[Information content]
The product $\Gamma(\tfrac14)\cdot\Gamma(\tfrac34)=\pi\sqrt{2}$ maps two
algebraically independent transcendentals to a single value algebraic in~$\pi$.
The individual values are not recoverable from the product alone.
The ratio $\Gstar=\Gamma(\tfrac14)/\Gamma(\tfrac34)$ preserves the relationship:
given $\Gstar$ and the product, both Gamma values are individually recoverable.
\end{observation}

\begin{conjecture}[Measurement--superposition correspondence]
The product operation (information-destroying) corresponds to quantum measurement:
$P=\abs{\psi}^2=\psi\cdot\psi^*$, which annihilates the complex phase.
The ratio operation (information-preserving) corresponds to superposition:
$\psi=a+bi$, which maintains the relationship between real and imaginary parts.

\smallskip
\begin{center}
\begin{tabular}{@{}llll@{}}
\toprule
Operation & Information & Quantum mechanics & Interpretation \\
\midrule
Product (collapse) & Destroyed & $P=\abs{\psi}^2$ & Measurement \\
Ratio (distinction) & Preserved & $\psi=a+bi$ & Superposition \\
\bottomrule
\end{tabular}
\end{center}
\end{conjecture}

The Born rule is a product: multiply $\psi$ by $\psi^*$ and the phase vanishes.
Superposition is a ratio: the wave function maintains the relationship between its
components. The phase---the thing proportional to~$i$---is what distinguishes
$\psi$ from $\psi^*$.

%==========================================================================
\section{The Imaginary Unit as Observer}
%==========================================================================

The blind derivation chain \cite{Steinmetz2026b} begins from a single axiom:
$i$~exists (i.e., $x^2+1=0$ has a solution). From this:
\begin{enumerate}[nosep]
  \item $\ZZ[i]$ exists (Gaussian integers = square lattice).
  \item $E_i\colon y^2=x^3-x$ is the unique CM curve with $j=1728$.
  \item $\abs{\Aut(E_i)}=4$, giving coefficient $16=4^2$.
  \item $\Gstar=\Gamma(\tfrac14)/\Gamma(\tfrac34)$ from the periods of~$E_i$.
  \item The master quadratic~\eqref{eq:quadratic} yields $x_+=137.036$.
  \item A one-loop lattice correction on $\ZZ[i]^3$ gives $x_+=137.036\,000$
        (9.6\,ppb from NIST).
\end{enumerate}
The entire structure of the coupling constants follows from~$i$.

\begin{conjecture}[The observer as kernel]\label{conj:observer}
The projection $\re\colon\CC\to\RR$ is the existence filter. Its kernel
$\ker(\re)=i\RR$ is the space of quantities that participate in the dynamics
but are projected out upon measurement. In quantum mechanics, this is the
global phase of the wave function. We identify $\ker(\re)$ with the observer:
the direction that is always perpendicular to every observable, always present
in the evolution, and always annihilated when one asks for an outcome.
\end{conjecture}

%==========================================================================
\section{Inversion: Physics as Constraint}
%==========================================================================

If $i$ is the ontological primitive---and the derivation chain suggests it
is---then the relationship between mathematics and physics inverts:

\begin{quote}
\emph{Standard view:} Physics describes what exists. Mathematics is the language.

\smallskip
\emph{Proposed view:} $i$ describes what exists (the full complex plane,
including the observer). Physics describes what survives projection onto
the real axis---what can be measured, observed, communicated.
\end{quote}

The laws of physics, on this reading, are not prescriptions for what happens.
They are constraints on what can be seen. The fine structure constant
$\alpha^{-1}=137.036$ is not a property of the electron; it is a property of
the projection $\re\colon\CC\to\RR$ applied to the arithmetic of~$\ZZ[i]$.

\begin{conjecture}[Alpha as projection cost]
$\alpha^{-1}=x_+$ is the cost of collapsing the complex plane onto the real
line at the scale set by the CM elliptic curve~$E_i$.
\end{conjecture}

%==========================================================================
\section{Why the Ratio Was Overlooked}
%==========================================================================

Four historical factors explain the omission:
\begin{enumerate}[nosep]
  \item \textbf{Period tradition.} Euler, Gauss, Legendre, and Jacobi studied
    the lemniscate and obtained $\varpi=\Gamma(\tfrac14)^2/(2\sqrt{2\pi})$.
    This involves the product, not the ratio.
  \item \textbf{$\pi$-centrism.} Writing constants in terms of~$\pi$ makes the
    product natural ($\pi\sqrt{2}$) and the ratio invisible ($\Gstar$ has no
    $\pi$-based expression).
  \item \textbf{Solved-constant bias.} The product reduces to~$\pi$ (closed form).
    The ratio does not. Mathematics gravitates toward solvable quantities.
  \item \textbf{Trivial period ratio.} For $E_i$, the period ratio
    $\omega_2/\omega_1=i$ is a known CM point; nobody extracts $\Gstar$ from it
    because the ratio of \emph{periods} is trivially~$i$---the ratio of
    \emph{Gamma values} is not.
\end{enumerate}

%==========================================================================
\section{Epistemic Status}
%==========================================================================

\begin{center}
\begin{tabular}{@{}lll@{}}
\toprule
Claim & Status & Basis \\
\midrule
Reflection formula produces product and ratio & Theorem & Algebraic identity \\
$\Gstar$ algebraically independent of $\pi$ & Theorem & Nesterenko 1996 \\
$\log\Gstar$ absorbs all unsolved $L$-values & Theorem & Identity~\eqref{eq:loggstar} \\
Product = collapse, ratio = distinction & Conjecture & Structural analogy \\
Observer $=\ker(\re)=i\RR$ & Conjecture & Interpretive \\
$\alpha^{-1}$ = projection cost & Conjecture & Interpretive \\
\bottomrule
\end{tabular}
\end{center}

\smallskip
The mathematical content (Sections~1--3) is rigorous. The interpretive content
(Sections~4--6) is conjecture. No experiment proposed here would distinguish this
interpretation from standard quantum mechanics. The value is conceptual: it provides
a structural reason \emph{why} $\alpha$ has the value it does and \emph{why} the
observer is absent from standard physics.

\bigskip
\noindent\textbf{Acknowledgments.}
The author thanks the developers of mpmath \cite{mpmath} for the arbitrary-precision
arithmetic library used throughout.

\begin{thebibliography}{99}

\bibitem{Chudnovsky1980}
G.~V.~Chudnovsky,
\emph{Contributions to the Theory of Transcendental Numbers},
AMS, 1984;
C.~R.~Acad.~Sci.~Paris \textbf{288}, A607 (1979).

\bibitem{Nesterenko1996}
Yu.~V.~Nesterenko,
Sb.~Math.\ \textbf{187}, 1319 (1996).

\bibitem{Steinmetz2026a}
W.~J.~Steinmetz~III,
``A quadratic in $\Gamma(\tfrac14)/\Gamma(\tfrac34)$ whose roots approximate
$\alpha^{-1}$ and $\alpha_s^{-1}$,''
preprint (2026).

\bibitem{Steinmetz2026b}
W.~J.~Steinmetz~III,
``The blind derivation: from $i$ to $\alpha^{-1}$ in 13 steps,''
FTD internal document FOUND\_BLIND\_DERIVATION\_CHAIN.md (2026).

\bibitem{mpmath}
F.~Johansson,
mpmath: a Python library for arbitrary-precision floating-point arithmetic (v1.3),
\url{https://mpmath.org} (2023).

\end{thebibliography}

\end{document}
